%! Linear Transformations — Unit 5, Lesson 5.3 — Guided Notes (Answer Key)
\documentclass[10pt]{article}
\usepackage{linalg-article}
\usepackage{linalg-key}   % pulls in -boxes; do NOT also load -boxes
\newcommand{\TallMath}[1]{$\displaystyle #1\rule[-1.4em]{0pt}{3.2em}$}
\newcommand{\vv}{\mathbf{v}}
\newcommand{\ww}{\mathbf{w}}
\newcommand{\xx}{\mathbf{x}}
\newcommand{\yy}{\mathbf{y}}
\newcommand{\ee}{\mathbf{e}}
\newcommand{\T}{\mathsf{T}}
\newcommand{\zero}{\mathbf{0}}

\begin{document}

\pageheader{Unit 5, Lesson 5.3}{Guided Notes --- Answer Key}
\namedateperiod

\vspace{0.12in}

\begin{objectivebox}
\textbf{By the end of this lesson, I will be able to\dots}
\begin{itemize}\setlength{\itemsep}{2pt}
  \item read a matrix as a \emph{motion of the plane} $\xx\mapsto A\xx$, and find where the unit square lands (the parallelogram of the columns);
  \item name the motion --- \emph{stretch, rotation, reflection, shear} --- from the matrix;
  \item use $|\det A|$ as the factor \emph{every} area is multiplied by, and read the \emph{sign} of $\det A$ (orientation) and $\det A=0$ (collapse $\Rightarrow$ no inverse).
\end{itemize}
\end{objectivebox}

\vspace{0.06in}

\begin{vocabbox}
\small
Fill in each term as we build it together.
\vspace{2pt}
\vocabans{Linear transformation ($\xx\mapsto A\xx$)}{a matrix read as a motion of the plane; it fixes the origin and keeps grid lines straight and evenly spaced}
\vocabans{Images of the basis ($A\ee_1,A\ee_2=$ columns)}{the columns of $A$ are where $\ee_1,\ee_2$ land, so the unit square maps to the column parallelogram}
\vocabans{Area (volume) factor $=|\det A|$}{applying $A$ multiplies every area (and every 3-D volume) by $|\det A|$}
\vocabans{Orientation (sign of $\det A$)}{$\det A>0$ preserves orientation, $\det A<0$ flips it; $\det A=0$ collapses space flat (no inverse)}
\end{vocabbox}

\begin{hookbox}
Photo and animation software rotates, stretches, and mirror-flips images by multiplying every pixel's
position vector by a $2\times2$ \textbf{matrix} --- the matrix moves the whole plane at once.
\begin{itemize}\setlength{\itemsep}{1pt}
  \item If I know the matrix, what one number tells me how much bigger the image gets? \ansline{The determinant --- $|\det A|$ is the factor every area is multiplied by.}
  \item And what tells me whether the picture came out mirror-flipped? \ansline{The \emph{sign} of the determinant: $\det A<0$ means orientation was flipped (a reflection).}
\end{itemize}
\end{hookbox}

\begin{notesbox}{1. A Matrix Moves Space}
Reading $\xx\mapsto A\xx$ as a \textbf{transformation}: every point of the plane is sent to a new point.
Because $A(\xx+\yy)=A\xx+A\yy$ and $A(c\xx)=cA\xx$, the origin stays fixed and straight, evenly spaced
grid lines stay straight and evenly spaced --- so the unit square maps to a \textbf{parallelogram}. And
$A\ee_1$, $A\ee_2$ are exactly the \emph{columns} of $A$ (Lesson~1.3), so watching the square tells you
the whole map. Here is $A=\begin{bsmallmatrix} 3 & 1 \\ 0 & 2 \end{bsmallmatrix}$:

\begin{center}
\begin{tikzpicture}[scale=0.62]
  % left: unit square
  \draw[linegray, thin, ->] (-0.3,0)--(1.8,0);
  \draw[linegray, thin, ->] (0,-0.3)--(0,2.0);
  \fill[burgundy!10] (0,0)--(1,0)--(1,1)--(0,1)--cycle;
  \draw[charcoal, thin] (0,0)--(1,0)--(1,1)--(0,1)--cycle;
  \draw[->, very thick, burgundy] (0,0)--(1,0);
  \draw[->, very thick, royalblue] (0,0)--(0,1);
  \node[charcoal, font=\scriptsize] at (0.5,0.5){area $1$};
  \node[burgundy, font=\scriptsize, below right] at (0.7,0){$\ee_1$};
  \node[royalblue, font=\scriptsize, left] at (0,0.7){$\ee_2$};
  % arrow
  \draw[->, very thick, charcoal] (2.2,0.9)--(3.2,0.9) node[midway,above,font=\scriptsize]{$A$};
  % right: image parallelogram, shifted
  \begin{scope}[shift={(3.7,0)}]
    \draw[linegray, thin, ->] (-0.3,0)--(4.6,0);
    \draw[linegray, thin, ->] (0,-0.3)--(0,2.4);
    \fill[royalblue!9] (0,0)--(3,0)--(4,2)--(1,2)--cycle;
    \draw[charcoal, thin] (0,0)--(3,0)--(4,2)--(1,2)--cycle;
    \draw[->, very thick, burgundy] (0,0)--(3,0);
    \draw[->, very thick, royalblue] (0,0)--(1,2);
    \node[charcoal, font=\scriptsize] at (2.05,0.95){area $|\det A|=6$};
    \node[burgundy, font=\scriptsize, below] at (3,0){$(3,0)$};
    \node[royalblue, font=\scriptsize, above left] at (1,2){$(1,2)$};
  \end{scope}
\end{tikzpicture}
\end{center}

\vspace{-2pt}
\textbf{Check the columns.} $A\ee_1=A\begin{bsmallmatrix} 1\\ 0\end{bsmallmatrix}=\begin{bsmallmatrix}
{\color{keyred}\mathbf{3}}\\ {\color{keyred}\mathbf{0}}\end{bsmallmatrix}$ and $A\ee_2=A\begin{bsmallmatrix} 0\\ 1\end{bsmallmatrix}
=\begin{bsmallmatrix} {\color{keyred}\mathbf{1}}\\ {\color{keyred}\mathbf{2}}\end{bsmallmatrix}$ --- the columns of $A$. A square
becomes a \ans{parallelogram}.
\end{notesbox}

\begin{notesbox}{2. A Gallery of Motions}
Read each matrix as a motion of the plane. (Recall a $2\times2$ determinant is $ad-bc$.)
\begin{center}
\renewcommand{\arraystretch}{1.5}
\begin{tabular}{@{}c l c@{}}
\textbf{Matrix} & \textbf{Motion} & $\det$ \\
\hline
$\begin{bsmallmatrix} 2 & 0 \\ 0 & 3 \end{bsmallmatrix}$ & stretch: $\times2$ across, $\times3$ up & $6$ \\
$\begin{bsmallmatrix} 0 & -1 \\ 1 & 0 \end{bsmallmatrix}$ & rotate $90^\circ$ (counterclockwise) & {\color{keyred}\textbf{1}} \\
$\begin{bsmallmatrix} 1 & 0 \\ 0 & -1 \end{bsmallmatrix}$ & reflect across the $x$-axis & {\color{keyred}$\mathbf{-1}$} \\
$\begin{bsmallmatrix} 1 & 1 \\ 0 & 1 \end{bsmallmatrix}$ & shear (slide the top sideways) & $1$ \\
\end{tabular}
\end{center}
Notice the rotation, reflection, and shear all have $|\det|=1$ --- they \emph{move} the square without
changing its \emph{area}.
\end{notesbox}

\begin{notesbox}{3. The Determinant Is the Area Factor}
The unit square has area $1$, and it lands on the parallelogram spanned by the columns --- whose area is
$|\det A|$. So applying $A$ multiplies \textbf{every} area by the same factor:
\[
  \boxed{\;\text{(new area)} = |\det A|\cdot(\text{old area})\;}
\]
\textbf{Worked example.} For $A=\begin{bsmallmatrix} 3 & 1 \\ 0 & 2 \end{bsmallmatrix}$ we have
$\det A=6$, so the unit square (area $1$) becomes a parallelogram of area \ans{$6$}. A triangle of
area $4$ becomes a shape of area $|\det A|\cdot 4={\color{keyred}\mathbf{24}}$. Every region is scaled by the same $6$.
\end{notesbox}

\begin{notesbox}{4. Sign, Collapse, and the Road Ahead}
The \emph{size} of $\det A$ is the area factor; its \emph{sign} carries one more piece of information.
\begin{itemize}\setlength{\itemsep}{3pt}
  \item \textbf{Orientation.} $\det A>0$ keeps orientation; $\det A<0$ \textbf{flips} it (a reflection is
        in the mix). A reflection has $\det=$ \ans{$-1$}: it keeps area (factor $1$) but reverses the picture.
  \item \textbf{Collapse.} If the columns are parallel, the image is \emph{flat} --- a line, area $0$. For
        $\begin{bsmallmatrix} 1 & 2 \\ 2 & 4 \end{bsmallmatrix}$, $\det=$ \ans{$0$}, so the unit square
        is squashed onto a line. A flattened plane cannot be un-flattened, so $\det A=0\Leftrightarrow A$ has no \ans{inverse}.
  \item \textbf{Composition.} Doing $B$ then $A$ is the matrix $AB$, and the area factors \emph{multiply}:
        $|\det(AB)|=|\det A|\,|\det B|$ --- the product rule from Lesson~5.2, now seen as area.
\end{itemize}
\vspace{2pt}
\textbf{The road ahead.} \textbf{Unit~6:} under a transformation most vectors get rotated, but a few
special directions are only \emph{stretched} (never turned). Those directions --- \emph{eigenvectors} ---
are the key to how a transformation behaves when you apply it again and \ans{again (repeatedly)}.
\end{notesbox}

\begin{practicebox}
Show your work. Use ``the columns are where $\ee_1,\ee_2$ land'' and ``area factor $=|\det A|$.''
\begin{enumerate}[leftmargin=1.6em, itemsep=6pt]
  \item For $A=\begin{bsmallmatrix} 2 & 0 \\ 0 & 3 \end{bsmallmatrix}$: where do $\ee_1$ and $\ee_2$ land, what is $\det A$, and what is the area of the image of a rectangle with area $5$? \ansline{$A\ee_1=(2,0)$, $A\ee_2=(0,3)$ (the columns). $\det A=6$. Area factor $|\det A|=6$, so the image has area $6\cdot5=30$.}
  \item Which flips orientation --- a reflection $\begin{bsmallmatrix} 1 & 0 \\ 0 & -1 \end{bsmallmatrix}$ or a stretch $\begin{bsmallmatrix} 2 & 0 \\ 0 & 2 \end{bsmallmatrix}$? How do you know from the determinant? \ansline{The reflection: $\det=-1<0$, so orientation flips. The stretch has $\det=4>0$, so orientation is preserved. The \emph{sign} of $\det$ tells you.}
  \item A matrix $A$ has $\det A=0$. What does its transformation do to the unit square, and is $A$ invertible? \ansline{It collapses the square flat --- onto a line or a point, area $0$. Since a flattened plane cannot be un-flattened, $A$ is \emph{not} invertible.}
\end{enumerate}
\end{practicebox}

\begin{teachernote}
The spine: a matrix \emph{is} a motion of the plane, and $|\det A|$ is the area factor of that motion.
Section~1 lands the picture --- $A\ee_1,A\ee_2$ are the columns, so the unit square becomes the column
parallelogram, whose area is $|\det A|=6$ (the Warm-Up made this concrete). Section~2 reads matrices as
motions and notices rotation/reflection/shear all preserve area. Section~3 is the payoff students will
use: new area $=|\det A|\cdot$ old area. Section~4 separates the two jobs of $\det$: size $=$ area
factor, sign $=$ orientation, and $\det=0=$ collapse $=$ ``singular, no inverse'' (the Unit~2 callback),
with composition giving the 5.2 product rule as multiplying area factors. Most common slips: dropping the
absolute value (an area is never negative), thinking a reflection shrinks area, and forgetting that
$\det=0$ means no inverse. The road-ahead line seeds Unit~6 (eigenvectors $=$ directions only stretched).
\end{teachernote}

\end{document}
